% Requiere en el preámbulo:
% \usepackage{booktabs}
% \usepackage{graphicx}
% \usepackage{tabularx}
% \usepackage{float}
% \usepackage{siunitx}
%
% Ajustar la ruta si esta carpeta no está junto al documento principal.

\section{Análisis descriptivo del corpus}
\label{sec:analisis-descriptivo}

Se realizó un análisis descriptivo del corpus MentalRiskES a partir de los
archivos JSON disponibles en las particiones \textit{Train}, \textit{Trial} y
\textit{Test}. Cada fila del conjunto de datos representa un mensaje y se
conservaron como metadatos el trastorno, la partición y el identificador del
sujeto. En total se analizaron 28\,085 mensajes pertenecientes a 820 sujetos.

\subsection{Composición del corpus}

La distribución de mensajes por trastorno fue la siguiente: 12\,066 mensajes
de depresión (42,96\%), 10\,499 de trastornos de la conducta alimentaria
(TCA; 37,38\%) y 5\,520 de ansiedad (19,66\%). Por partición, el conjunto de
\textit{Test} contiene 14\,893 mensajes (53,02\%), \textit{Train} 12\,179
(43,36\%) y \textit{Trial} 1\,013 (3,61\%).

\begin{table}[H]
    \centering
    \caption{Número de mensajes por trastorno y partición.}
    \label{tab:distribucion-trastorno-particion}
    \begin{tabular}{lrrr}
        \toprule
        Trastorno & Train & Trial & Test \\
        \midrule
        Ansiedad  & 0 & 0 & 5\,520 \\
        Depresión & 6\,248 & 624 & 5\,194 \\
        TCA       & 5\,931 & 389 & 4\,179 \\
        \midrule
        Total     & 12\,179 & 1\,013 & 14\,893 \\
        \bottomrule
    \end{tabular}
\end{table}

\begin{figure}[H]
    \centering
    \includegraphics[width=0.88\textwidth]{images/Analisis_descriptivo/distribucion_particion_trastorno.png}
    \caption{Mapa de calor del número de mensajes por partición y trastorno.}
    \label{fig:distribucion-trastorno-particion}
\end{figure}

\subsection{Distribución de las clases}
\label{sec:distribucion-clases}

La distribución de clases se calculó respetando el esquema de anotación
específico de cada trastorno. Para ansiedad y TCA se utilizó la etiqueta
binaria de \texttt{gold\_a.txt}, con las categorías \textit{control} y
\textit{suffer}. Para depresión se utilizó la variable clínica de
\texttt{gold\_c.txt}, que distingue entre \textit{control},
\textit{suffer+in favour}, \textit{suffer+against} y
\textit{suffer+other}. Los porcentajes de la Tabla~\ref{tab:clases-sujetos}
se calculan sobre sujetos, mientras que los de la
Tabla~\ref{tab:clases-mensajes} se calculan sobre mensajes.

\begin{table}[H]
    \centering
    \caption{Distribución de clases por sujeto y partición. Los valores se
    expresan como número de sujetos y porcentaje dentro de cada combinación
    de trastorno y partición.}
    \label{tab:clases-sujetos}
    \resizebox{\textwidth}{!}{%
    \begin{tabular}{llrrrrrr}
        \toprule
        Trastorno & Partición & Control & Suffer &
        Suffer+in favour & Suffer+against & Suffer+other & Total \\
        \midrule
        Ansiedad  & Test  & 57 (38,0\%) & 93 (62,0\%) & -- & -- & -- & 150 \\
        \midrule
        Depresión & Train & 81 (46,3\%) & -- & 44 (25,1\%) &
        44 (25,1\%) & 6 (3,4\%) & 175 \\
        Depresión & Trial & 4 (40,0\%) & -- & 2 (20,0\%) &
        2 (20,0\%) & 2 (20,0\%) & 10 \\
        Depresión & Test  & 81 (54,0\%) & -- & 32 (21,3\%) &
        32 (21,3\%) & 5 (3,3\%) & 150 \\
        \midrule
        TCA       & Train & 101 (57,7\%) & 74 (42,3\%) & -- & -- & -- & 175 \\
        TCA       & Trial & 5 (50,0\%) & 5 (50,0\%) & -- & -- & -- & 10 \\
        TCA       & Test  & 86 (57,3\%) & 64 (42,7\%) & -- & -- & -- & 150 \\
        \bottomrule
    \end{tabular}%
    }
\end{table}

\noindent
\footnotesize
\textit{Nota:} ``--'' indica que esa categoría no forma parte del esquema de
anotación del trastorno o que no aparece en la partición correspondiente. En
ansiedad y TCA las clases son binarias (\textit{control} y \textit{suffer}).
En depresión se emplean las cuatro categorías clínicas de
\texttt{gold\_c.txt}. La columna ``Total'' permite comprobar que los
recuentos de cada fila suman el número total de sujetos de la partición.
\normalsize

\begin{table}[H]
    \centering
    \caption{Distribución de clases por mensaje.}
    \label{tab:clases-mensajes}
    \scriptsize
    \begin{tabular}{llrrrrr}
        \toprule
        Trastorno & Partición & Control & Suffer & Suffer+in favour &
        Suffer+against & Suffer+other \\
        \midrule
        Ansiedad  & Test  & 2\,222 (40,3\%) & 3\,298 (59,7\%) & -- & -- & -- \\
        Depresión & Train & 3\,135 (50,2\%) & -- & 1\,524 (24,4\%) &
        1\,457 (23,3\%) & 132 (2,1\%) \\
        Depresión & Trial & 327 (52,4\%) & -- & 35 (5,6\%) &
        136 (21,8\%) & 126 (20,2\%) \\
        Depresión & Test  & 2\,825 (54,4\%) & -- & 1\,154 (22,2\%) &
        1\,072 (20,6\%) & 143 (2,8\%) \\
        TCA       & Train & 3\,399 (57,3\%) & 2\,532 (42,7\%) & -- & -- & -- \\
        TCA       & Trial & 228 (58,6\%) & 161 (41,4\%) & -- & -- & -- \\
        TCA       & Test  & 2\,959 (70,8\%) & 1\,220 (29,2\%) & -- & -- & -- \\
        \bottomrule
    \end{tabular}
\end{table}

\begin{figure}[H]
    \centering
    \includegraphics[width=0.90\textwidth]{images/Analisis_descriptivo/distribucion_clases.png}
    \caption{Composición porcentual de las clases anotadas por trastorno y partición.}
    \label{fig:distribucion-clases}
\end{figure}

La comparación entre las tablas por sujeto y por mensaje muestra que la
distribución de mensajes no coincide necesariamente con la distribución de
sujetos. Esto se debe a que los sujetos presentan un número variable de
mensajes. Por ejemplo, en el conjunto de \textit{Test} de TCA la clase
\textit{control} representa el 57,3\% de los sujetos, pero el 70,8\% de los
mensajes.

\subsection{Calidad de los datos}

No se detectaron mensajes nulos ni cadenas vacías. Se identificaron 31
mensajes duplicados exactos al comparar el contenido textual de los mensajes,
sin utilizar los identificadores de mensaje para dicha comparación. Estos
duplicados no se eliminaron, ya que pueden corresponder a repeticiones
legítimas en conversaciones distintas.

\subsection{Longitud de los mensajes}

La longitud se calculó como el número de palabras separadas por espacios. La
longitud media fue de 14,70 palabras y la mediana de 9 palabras, con una
desviación estándar de 26,38. El primer y tercer cuartil fueron 6 y 16
palabras, respectivamente, mientras que el mensaje más largo alcanzó 899
palabras. La diferencia entre la media y la mediana, junto con el máximo
observado, indica una distribución asimétrica con presencia de mensajes
considerablemente más largos.

\begin{table}[H]
    \centering
    \caption{Estadísticas descriptivas de la longitud de los mensajes.}
    \label{tab:longitud-mensajes}
    \begin{tabular}{lr}
        \toprule
        Estadístico & Palabras \\
        \midrule
        Número de mensajes & 28\,085 \\
        Media & 14,70 \\
        Desviación estándar & 26,38 \\
        Mínimo & 1 \\
        Primer cuartil (Q1) & 6 \\
        Mediana & 9 \\
        Tercer cuartil (Q3) & 16 \\
        Máximo & 899 \\
        \bottomrule
    \end{tabular}
\end{table}

Como representación gráfica de la longitud de los mensajes se utiliza un
diagrama de cajas. La Figura~\ref{fig:boxplot-longitud} muestra la mediana,
el rango intercuartílico y los valores atípicos de cada combinación de
trastorno y partición. Para mejorar la legibilidad, el eje vertical se limita
al percentil 99 de la longitud global; los mensajes que superan este valor no
se eliminan del análisis y sí se incluyen en las tablas descriptivas.

\begin{figure}[H]
    \centering
    \includegraphics[width=\textwidth]{images/Analisis_descriptivo/boxplot_longitud_por_grupo.png}
    \caption{Diagrama de cajas de la longitud de los mensajes por trastorno y partición.}
    \label{fig:boxplot-longitud}
\end{figure}

El boxplot permite observar que la distribución de la longitud es asimétrica
en todos los grupos: las medianas se sitúan entre 8 y 10 palabras, mientras
que existe una cola de mensajes más extensos. Esta variabilidad es
especialmente relevante para el preprocesamiento, ya que una representación
basada únicamente en la longitud media podría ocultar la diferencia entre
mensajes breves y mensajes largos con mayor contenido discursivo.

Se detectaron 122 mensajes de tres palabras o menos. De ellos, 34 (27,87\%)
fueron clasificados heurísticamente como emoticonos o emojis, mediante
patrones textuales y descripciones frecuentes de emojis. Esta identificación
es aproximada y no debe interpretarse como una anotación manual.

\subsection{Vocabulario}

La tokenización básica produjo 382\,124 tokens y un vocabulario de 26\,785
palabras únicas. De estas, 14\,810 aparecen una sola vez
(\textit{hapax legomena}), lo que representa el 55,29\% del vocabulario. Las
palabras más frecuentes fueron ``de'', ``que'', ``y'', ``no'', ``la'', ``a'',
``me'', ``es'', ``en'' y ``el''. La elevada proporción de palabras con
frecuencia uno sugiere una notable diversidad léxica y justifica evaluar
normalización, filtrado de ruido y tratamiento de variantes ortográficas antes
de entrenar modelos de representación textual.

\begin{figure}[H]
    \centering
    \includegraphics[width=0.82\textwidth]{images/Analisis_descriptivo/palabras_mas_frecuentes.png}
    \caption{Quince palabras más frecuentes del corpus, representadas mediante un gráfico de tipo \textit{lollipop}.}
    \label{fig:palabras-frecuentes}
\end{figure}
